% =====================================================================
% 第 6 章  讨论与结论
% ---------------------------------------------------------------------
% 编译说明同前：ctexart 下用 \section，ctexbook/report 下改为 \chapter。
% =====================================================================

\section{讨论与结论}

\subsection{主要发现}

本文用 2001 年至 2026 年的 102 期《货币政策执行报告》检验了文本情绪的定价效应，
得到三点结果。

第一，政策指引段的情绪对发布后的市场收益方向没有稳健影响。全样本中，一日与五日
窗口的系数都接近于零，其标准误对应的可识别区间在数十个基点的量级，说明这一零结果
不是估计失败，而是本样本确实识别不到方向性效应。

第二，同一情绪显著降低了发布日的市场波动。无论以绝对收益还是日内振幅度量，控制
事前波动之后系数仍为负且在 5\% 水平显著，28 个行业的口径给出几乎相同的估计，
置换检验的 $p$ 值为 0.020。

第三，把样本按 2014 年切开后，波动效应集中在前期，方向效应出现在后期且为负。
这一转折对断点年份的选择不敏感，却对剔除 2015 年较为敏感，因此本文把它作为提示性
证据，而非与第二条并列的结论。

\subsection{为什么作用在波动上而不是方向上}

这一结果的直觉并不复杂。政策文件的作用首先是让市场对政策路径形成更一致的判断，
而不是直接告诉市场资产应该值多少钱。当报告传递的指引较为明确时，投资者对未来的
分歧收敛，交易行为随之趋同，价格的单日摆动幅度下降；至于价格往哪个方向走，还取决于
基本面、流动性与估值，单一文件的措辞难以决定。

这一解释与既有证据是相容的。姜富伟等（2021）在四类宽基指数上发现，文本情绪上升
伴随收益上升，文本相似度上升伴随波动下降；两者的共同点在于，文本提供的是关于政策
路径的信息，其对分歧的作用比对水平的作用更直接。王熙和郑梦圆（2021）
把语调指数纳入结构向量自回归，控制信息效应之后“价格之谜”消失，进一步说明央行文本
中确实含有关于未来政策的信息成分，而不只是对已发生事实的复述。

本文的分段结果为这一解释提供了额外支持。若波动下降只是“整篇报告听起来更乐观”的
结果，那么宏观回顾段的情绪应当具有类似作用；实际估计中，回顾段的系数为 $-2.1$
个基点（$t=-0.20$），几乎为零。这提示起作用的是面向未来的部分，与姜富伟等（2021）
关于“政策指引而非宏观回顾”的判断一致。王辉等（2023）把报告拆为前瞻性与回顾性两部分
分别度量，蒋海等（2025）则以机器学习方法提取前瞻性情绪，两项研究的设计都建立在
“前瞻内容承载主要政策信号”这一前提之上。

\subsection{2014 年之后的方向性变化}

方向性效应在 2014 年之后才出现，且符号为负。本文认为有两种互不排斥的解释。

第一种是“坏消息”效应。宽松措辞往往出现在经济下行压力较大的时点，措辞越积极，
越可能意味着政策需要发力的程度越深。市场若把宽松解读为对基本面的确认，报告发布后
反而出现负收益。这一机制与王宇伟等（2025）关于央行言辞沟通非对称性的发现相通：
在他们关于企业家信心的研究中，紧缩表述有效、宽松表述效果有限，说明市场对不同方向
的措辞并非对称处理。

第二种是预期管理的成熟。2014 年之后，央行在报告发布前的沟通渠道明显增多，
政策意图被提前定价，等到报告正式发布，新增信息有限，“利好出尽”式的回调更容易出现。
这一解释可以说明为什么方向效应为负而波动效应同时减弱：当信息已被消化，文件本身
既不能推动价格，也难以收敛分歧。

需要承认的是，本文的数据无法在两种解释之间作出裁决。可用的证据只有一条：剔除 2015
年后方向效应明显减弱（交互项由 $-75.6$ 降至 $-59.0$，$t=-1.62$），剔除 2008--2009 年
后进一步降至 $-45.5$（$t=-1.56$），两者都不再显著。这说明该方向效应相当一部分由金融
危机与股市异常波动期的数期报告贡献，而这类时期恰好是“坏消息”效应最可能显现的时点。
这一点更倾向于第一种解释，但不足以排除第二种。

\subsection{与既有研究的关系}

与姜富伟等（2021）相比，本文的差别有两处。一是样本延长至 2026 年，覆盖了 2020 年
之后的非常规货币政策阶段；二是把被解释变量从收益扩展到波动，并引入制度断点。
结论上，本文确认了“指引段而非回顾段”的判断，但没有复现出显著的方向效应。这一差异
可能来自三方面：分段依据不同（本文按章节，原文按句子内容）、市场代理不同（本文用
行业等权组合，原文用四类宽基指数）、样本期不同。第 3 章已就第一点作了说明。

与周阳、綦岚鑫（2024）相比，两者的差异主要在结果变量。他们构建了专门适配报告的
领域词典，并同时考察股票、债券与外汇市场的收益与波动，发现市场之间存在溢出，
溢出会放大或削弱情绪效应。本文只做股票市场，但在波动维度上得到的结论与之相容：
情绪对二阶矩的作用比对一阶矩更稳定。他们词典中纳入否定词与程度副词的做法，
也为本文后续改进度量提供了方向。

与游庭玮（2026）相比，两者都使用申万一级行业数据，但解释变量不同。他用美联储政策
利率，本文用报告文本；他通过控制市场整体收益识别行业独立冲击，本文则把行业等权组合
直接作为市场代理。两条路线并不冲突，若能取得覆盖期更长的行业特征数据，把文本情绪
与行业暴露度结合，是更自然的下一步。

陈妍玲等（2024）在金融业及其子行业上发现，前瞻性报告的文本情绪显著提高风险溢出，
这与本文“情绪影响二阶矩”的结论方向一致，只是他们关注的是跨市场溢出，本文关注的是
发布日的波动。林建浩和孙乐轩（2025）在综述中提醒，文本度量对模型设定与提示方式较为
敏感，可复现性有待改善。本文以两套词典与置换检验回应了这一提醒，也确实观察到替代
词典下系数衰减的现象。

\subsection{政策含义}

如果沟通的主要作用是收敛分歧，那么它的价值就不应只用“能否拉动指数”来衡量。
在预期管理框架下，一份表述清晰的报告可以降低市场对政策路径的不确定性，减少无谓的
波动，这本身即是政策目标的一部分。本文的估计为这一判断提供了数量上的支持：
情绪每提高一个标准差，发布日的日内振幅下降约 28 个基点，控制事前波动后仍有约
19 个基点。

另一方面，方向性效应的不稳定提示不宜高估单份文件的短期市场影响。2014 年之后方向
效应转为负，且对个别年份敏感，说明市场对报告的解读依赖于当时的宏观环境与预期状态。
对政策制定者而言，这意味着沟通的效果不是恒定的，需要结合所处的政策周期来判断；
对投资者而言，把报告措辞当作方向性交易信号，其可靠性远低于把它当作不确定性指标。

\subsection{局限与后续研究}

本文的局限主要有四。

其一，度量误差。词典法只能识别词面极性，对条件句、反讽与语境依赖的处理能力有限。
替代词典下系数按覆盖率比例收缩，说明误差主要是衰减型的，因而本文报告的效应量级
更可能是下界而非上界。引入带否定词与程度副词的领域词典，或用预训练模型替代词典，
是可行的改进方向。

其二，识别的时间精度。元数据不含发布的具体时刻，事件窗的起点只能以交易日为单位
近似。虽然更换窗口的结果一致，但对盘中反应与隔夜反应的区分仍无法做到。

其三，同期信息的干扰。货币政策操作、金融统计与物价数据常在相近时点公布，事件窗内
可能混入其他信息。本文的处理有两条：一是把降准、LPR 与逆回购操作作为控制变量加入方程，
加入后情绪系数几乎不变；二是把事件日平移若干交易日重做，结果发现效应在 $\pm$10 天以外
消失，但前移 5 天仍有约九成幅度，说明它并非严格的单日冲击。这两点缓解了担忧，但没有
完全消除。

其四，样本与功效。99 个事件、日频数据，能够可靠识别的效应量级在 70 个基点以上，
更细微的效应无法区分。若要把方向效应估计得更精确，需要把事件扩展到其他沟通文本，
例如例会纪要、新闻发布会实录，或改用更高频的数据。

行业层面是本文有意未展开的部分。第 4 章已说明，报告级变量在季度面板中与时间固定
效应共线，要做行业异质性，必须依赖行业特征与情绪的交互，而这要求行业特征在样本期内
有足够变异。现有数据中的估值与财务指标覆盖期偏短，难以支撑稳定估计。若后续能补齐
这部分数据，“哪些行业对政策指引更敏感”仍是一个值得回答的问题。

\subsection{结论}

本文的出发点是一个简单的问题：央行在报告里说的话，市场听进去了吗？答案取决于问的是
哪个方面。就价格方向而言，本文没有找到稳健的证据；就市场分歧而言，证据是清楚的——
政策指引段的措辞越积极，报告发布当日的波动越小，且这一作用并非来自报告的整体语气。

这一区分对理解央行沟通有意义，对评估沟通效果同样有影响。若只估计全样本的平均效应，
2014 年以前的方向性零结果与之后的方向性负效应会相互抵消，得到的结论是“沟通对价格
没有影响”；而真实情况是，沟通的作用先体现在波动上，随时间推移才在方向上有所显现。
把不同时期、不同文本段落视为同质样本，容易掩盖这一结构变化。
